% Nouveau EGA V (cote 25): archive cote 25, site page 43 (PDF page 44).
% typescript. Draft compared: none.
% AI-transcribed by Claude Opus 5.5 (claude-opus-5-5, adaptive thinking, effort high), Message Batches API, batch ega25b (msgbatch_01A7n5YuXubUPNK1zsepi8RG), 2026-11-13; not yet checked by eye.
\scanpage{25/43}
\mspage{\unsure{39}}

L'adhérence de son images dans $\overset{\hand{\vee}}{P}$ , muni de la structure réduite
induite, et considérons le morphisme induit g: $\underline{Y}^{\text{sing}} \longrightarrow T$ \xout{de} (morphis-%
me dominant de préschémas \xout{irréductibles} intègres). Alors les condition
\inserted[typed]{(i) à (ii bis)} suivantes sont équivalentes:

(i) Le morphisme g est génériquement \deleted{non ramifié} \inserted{étale} (i.e. est \deleted{non
ramifié} \inserted{étale} en au moins un point, ou ce qui revient au même, est \inserted{étale =} non ramifié
au point générique de $\underline{Y}^{\text{sing}}$).

(ibis) L'extension de corps L/K définie par g est finie et
séparable.

\margin{NB. L'hyp. \unsure{(on fera ici en} \illegible{1} \unsure{lien de : non normalité)} ne sert que pour l'équivalence avec (i ter). Énoncer 8.8. en conséquence\dots}
(i ter) Le morphisme g est birationnel, i.e. l'extension L/K est
l'extension triviale.

(ii) $\underline{Y}^{\text{sihg}} \neq \underline{Y}^{\text{sup sing}}$ \inserted{(ensemblistement, disons).}

(ii bis) Il existe un $x \in X(\bar{k})$ , et un \over{h}{h}yperplan tangent \inserted{H} à $X_{\bar{k}}$
en x qui ne soit pas ``osculateur'' en x (par quoi on entenda précisément
que x n'est pas supersingulier pour la section de $\underline{O}_{X_{\bar{k}}}(1)$ qui définit
H ... ).

Ces conditions impliquent que \inserted{$\underline{Y}^{\text{sup sing}} \neq \emptyset$ \unsure{ou}} dim $\underline{Y}^{\text{sup sing}}$ \deleted{$\leq$}$= r-2$, donc que
l'image de $\underline{Y}^{\text{sup sing}}$ dans $\overset{\hand{\vee}}{P}$ a une \unsure{dodimension|codimension} $\geqslant 2$, et \xout{que T
est de dimension r-1, i.e. de codimension 1 dans P} elles impliquent
également:

(iii) dim T = r-1, i.e. T est de codimension 1 dans $\overset{\vee}{P}$ .

Dém. L'équivalence de (i) et (i bis) est triviale , son équiva-%
lence avec (ii) est une conséquence triviale de 8.\unsure{7}. b), enfin l'équi-%
valence de (ii) et (ii bis) est pratiquement la définition de
$\underline{Y}^{\text{sup sing}}$ . Evidemment (i ter) \hand{$\Longrightarrow$}(i) , il reste à prouver que (i)\hand{$\Longrightarrow$}
(i ter). On peut évidemment supposer k alg clos, et on est ramené à
prouver (compte tenu de l'hypothèse (\unsure{iii}) ) qu\xout{e}'il existe un ouvert U\hand{$\neq\emptyset$}

\margin{43}
